

اینترنت را اگر بخواهیم ساده و خودمانی توضیح بدهیم، می‌شود آن را به یک شهر خیلی بزرگ شبیه کرد؛ شهری که در آن خانه‌ها، خیابان‌ها و چراغ‌راهنماها همه با هم کار می‌کنند تا رفت‌وآمدها درست و منظم باشد. در این شهر، «هاست»ها یا همان پایانه‌ها مثل خانه‌ها هستند. هاست همان لپ‌تاپ، گوشی یا سرور ماست که برنامه‌ها و اپلیکیشن‌ها روی آن اجرا می‌شوند. یعنی جایی است که پیام ساخته می‌شود یا پیام‌های دریافتی به آن می‌رسند.  

وقتی این پیام‌ها قرار است از یک خانه به خانه‌ی دیگر برسند، لازم است که کسی مسیر را نشان بدهد. این‌جاست که «سوییچ بسته»‌ها وارد عمل می‌شوند. سوییچ بسته‌ها مثل چراغ‌راهنما یا پلیس‌های چهارراه هستند که به هر بسته‌ی اطلاعاتی می‌گویند از کدام مسیر حرکت کند. این سوییچ‌ها دو مدل اصلی دارند: سوییچ‌های معمولی و «روتر»ها. سوییچ‌ها بیشتر در شبکه‌های محلی استفاده می‌شوند، ولی روترها نقش کلیدی‌تری دارند و شبکه‌های مختلف را به هم وصل می‌کنند.  

مسیرهایی هم که این بسته‌ها طی می‌کنند، گوناگون است. گاهی از کابل مسی رد می‌شوند، گاهی از فیبر نوری و گاهی هم بی‌سیم، مثل امواج رادیویی. فیبر نوری به خاطر سرعت و پهنای باند بالایش خیلی مهم است و امروز بخش بزرگی از اینترنت پرسرعت روی همین فیبر نوری کار می‌کند.  

اگر از دید سخت‌افزاری یا همان پیچ و مهره‌ای نگاه کنیم، اینترنت چیزی جز مجموعه‌ای از هاست‌ها، روترها و لینک‌های ارتباطی نیست. اما فقط این‌ها کافی نیستند، چون باید یک‌سری قوانین هم وجود داشته باشد تا همه بتوانند با هم هماهنگ شوند. این قوانین همان «پروتکل»ها هستند. پروتکل‌ها تعیین می‌کنند پیام چه شکلی داشته باشد، چطور فرستاده شود و وقتی به مقصد رسید، چه واکنشی به آن نشان داده شود.  

از یک دید دیگر، اینترنت بیشتر به چشم یک «سرویس» دیده می‌شود. در این دیدگاه، اینترنت یک زیرساخت بزرگ است که برای اپلیکیشن‌ها خدمات فراهم می‌کند. برای اینکه اپلیکیشن‌ها بتوانند از اینترنت استفاده کنند، یک رابط لازم است به نام \lr{API} که داده‌ها را آماده می‌کند و تحویل شبکه می‌دهد تا به مقصد برسند.  

«لبه‌ی شبکه» جایی است که کاربران به اینترنت وصل می‌شوند. گوشی، لپ‌تاپ یا سرور همه از طریق همین لبه وارد شبکه می‌شوند. اتصال می‌تواند با سیم باشد، مثل کابل اترنت، یا بی‌سیم، مثل وای‌فای و اینترنت موبایل.  

یک مثال ساده از شبکه‌ی خانگی این است که چند دستگاه مختلف در خانه از طریق وای‌فای به روتر وصل می‌شوند. بعضی وقت‌ها هم دستگاه‌ها مستقیم با کابل به روتر وصل می‌شوند که سرعت بیشتری دارد. روتر در خانه به مودم وصل است و مودم هم سیگنال‌های دریافتی از شرکت اینترنت را به شکل قابل فهم برای روتر ترجمه می‌کند.  

شبکه‌ها فقط خانگی نیستند؛ شبکه‌های بزرگ سازمانی یا دانشگاهی هم وجود دارند که ترکیبی از ارتباطات با سیم و بی‌سیم‌اند و پر از روتر و سوییچ هستند تا همه‌چیز منظم پیش برود.  

وقتی یک دستگاه می‌خواهد پیام بفرستد، آن را به بخش‌های کوچک‌تری به نام «پکت» تقسیم می‌کند و این پکت‌ها را از طریق شبکه‌ی دسترسی وارد اینترنت می‌کند. این پکت‌ها برای حرکت نیاز به یک بستر دارند که به آن «مدیا» می‌گوییم. مدیا دو نوع اصلی دارد: هدایت‌شده و غیرهدایت‌شده.  

مدیاهای هدایت‌شده مثل کابل مسی، کابل کواکسیال یا فیبر نوری‌اند. فیبر نوری داده را به شکل نور منتقل می‌کند، سرعت خیلی بالاست و چون نویز نمی‌گیرد، خطایش کم است. کابل کواکسیال هم دوطرفه است و می‌تواند چند کانال مختلف را هم‌زمان پشتیبانی کند.  
مدیاهای غیرهدایت‌شده همان امواج رادیویی‌اند که دیگر نیازی به کابل ندارند و داده را از طریق هوا منتقل می‌کنند. البته این‌ها ممکن است تحت تأثیر بازتاب یا برخورد به اجسام کیفیتشان پایین بیاید.  

روش‌های دیگری هم هست، مثل ارتباطات ماهواره‌ای یا مایکروویو زمینی که برای فاصله‌های خیلی زیاد یا جاهایی که کابل‌کشی سخت است، استفاده می‌شوند.  

در نهایت، اینترنت ترکیبی است از دستگاه‌ها (مثل گوشی و سرور)، مسیرها (مثل روتر و سوییچ) و رسانه‌های انتقال (مثل سیم و موج). پروتکل‌ها هم قوانین بازی هستند که همه‌چیز را منظم نگه می‌دارند. به همین دلیل است که ما می‌توانیم راحت در خانه بنشینیم و فیلم ببینیم یا پیامی که فرستاده‌ایم در چند ثانیه به آن سر دنیا برسد.  
